\documentclass[a4paper,12pt]{article}

\usepackage[brazil]{babel}
\usepackage[utf8]{inputenc}
\usepackage[T1]{fontenc}

\usepackage{geometry}
\usepackage{setspace}
\usepackage{titlesec}
\usepackage{hyperref}
\usepackage{graphicx}
\usepackage{caption}
\usepackage{subcaption}
\usepackage{fancyhdr}
\usepackage{xcolor}
\usepackage{amsmath}
\usepackage{amssymb}

%%%%%%%%%%%%%%%%%%%%%%%%%%%%%%%%%%%%%%%%%%%%%%%%%%
% These are some new commands that may be useful 
% for paper writing in general. If other new commands
% are needed for your specific paper, please feel 
% free to add here. 
%
% The currently available commands are organized in: 
% 1) Systems
% 2) Quantities
% 3) Energies and units
% 4) particle species
% 5) Colors package
% 6) hyperlink
%%%%%%%%%%%%%%%%%%%%%%%%%%%%%%%%%%%%%%%%%%%%%%%%%%

\usepackage{amsmath}
\usepackage{amssymb}
\usepackage{upgreek}
\usepackage{multirow}
\usepackage{setspace}% http://ctan.org/pkg/setspace
\usepackage{fancyhdr}
\usepackage{datetime}

% 1) SYSTEMS
\newcommand{\btc}               {\textbf{BTC}}
\newcommand{\btcspace}          {\textbf{BTC} }
\newcommand{\pow}               {\textbf{PoW}}

% 4) definition to references, biblatex and hyperlink
\usepackage[backend=bibtex, 
style=nature,  %style reference.
sorting=none,
firstinits=true %first name abbreviate
]{biblatex}

\usepackage{hyperref}
\hypersetup{
    colorlinks=true, %set "true" if you want colored links
    linktoc=all,     %set to "all" if you want both sections and subsections linked
    linkcolor=blue,  %choose some color if you want links to stand out
    citecolor= blue, % color of \cite{} in the text.
    urlcolor  = blue, % color of the link for the paper in references.
}

% 5) Tikz and figures
\usepackage{epsfig}
\usepackage{lmodern}
\usepackage{mathtools}
\usepackage[utf8]{luainputenc}
\usepackage{xspace}
\usepackage{tikz}
\usepackage{pgfplots}
\pgfplotsset{compat=newest}

\usetikzlibrary{positioning}
\usepackage{subcaption}

% 6) colors:
\usepackage{xcolor}
\definecolor{ao(english)}{rgb}{0.0, 0.5, 0.0} % dark green

% 7) Add lines numbers
%\usepackage{lineno}

% add pdf file to thesis:
\usepackage{pdfpages}

\usepackage[most]{tcolorbox}

%textmarker style from colorbox doc
\tcbset{textmarker/.style={%
        enhanced,
        parbox=false,boxrule=0mm,boxsep=0mm,arc=0mm,
        outer arc=0mm,left=6mm,right=3mm,top=7pt,bottom=7pt,
        toptitle=1mm,bottomtitle=1mm,oversize}}


% define new colorboxes
\newtcolorbox{hintBox}{textmarker,
    borderline west={6pt}{0pt}{yellow},
    colback=yellow!10!white}
\newtcolorbox{importantBox}{textmarker,
    borderline west={6pt}{0pt}{red},
    colback=red!10!white}
\newtcolorbox{noteBox}{textmarker,
    borderline west={6pt}{0pt}{green},
    colback=green!10!white}

% define commands for easy access
\newcommand{\note}[1]{\begin{noteBox} \textbf{} #1 \end{noteBox}}
\newcommand{\warning}[1]{\begin{hintBox} \textbf{Warning:} #1 \end{hintBox}}
\newcommand{\important}[1]{\begin{importantBox} \textbf{} #1 \end{importantBox}}

% --------------------------------------------------
% CORES E CAIXAS
% --------------------------------------------------

\definecolor{azul}{RGB}{30,90,150}
\definecolor{cinza}{RGB}{245,245,245}
\definecolor{verde}{RGB}{230,245,235}
\definecolor{amarelo}{RGB}{255,248,220}

\newtcolorbox{ideia}{
    colback=azul!5,
    colframe=azul,
    title=\textbf{IDEIA-CHAVE},
    fonttitle=\bfseries
}

\newtcolorbox{exemplo}{
    colback=verde,
    colframe=green!50!black,
    title=\textbf{EXEMPLO},
    fonttitle=\bfseries
}

\newtcolorbox{atencao}{
    colback=amarelo,
    colframe=orange!80!black,
    title=\textbf{ATENÇÃO},
    fonttitle=\bfseries
}

% --------------------------------------------------
% CABEÇALHO
% --------------------------------------------------

\pagestyle{fancy}
\fancyhf{}

\renewcommand{\headrulewidth}{0pt}

\fancyhead[R]{\thepage}

% --------------------------------------------------
% HYPERREF
% --------------------------------------------------

\hypersetup{
    colorlinks=true,
    linkcolor=blue,
    citecolor=blue,
    urlcolor=blue
}

% --------------------------------------------------
% BIBLIOGRAFIA
% --------------------------------------------------

\addbibresource{bibliography.bib}

\newcommand{\printingbibliography}{%
    \pagestyle{myheadings}
    \markright{}
    \sloppy
    \printbibliography[
        heading=bibintoc,
        title=Referências
    ]
    \fussy
}

% --------------------------------------------------
% CONFIGURAÇÕES
% --------------------------------------------------

\PassOptionsToPackage{table}{xcolor}

\renewcommand{\baselinestretch}{1.5}

\geometry{
    a4paper,
    top=30mm,
    bottom=20mm,
    left=30mm,
    right=20mm
}

\titleformat*{\section}{\bfseries\large}
\titleformat*{\subsection}{\bfseries\normalsize}

\title{
    \textbf{\large
    PROVÃO PAULISTA SERIADO 1ª, 2ª E 3ª SÉRIES DO ENSINO MÉDIO.}
}

\author{Andr\'e V. Silva \\ \href{https://andrevsilva.com/}{andrevsilva.com}}

\date{\today}


% ==================================================
% DOCUMENTO
% ==================================================

\begin{document}

\maketitle

\begin{center}
    \textbf{\Large CINEMÁTICA}\\[0.2cm]
    \textbf{\large 1.1 Velocidade e Aceleração}
\end{center}

\tableofcontents

\vspace{0.5cm}

\begin{ideia}

A \textbf{Cinemática} é a parte da Física que estuda e descreve
o movimento dos corpos.

Para compreender um movimento, podemos fazer perguntas como:

\begin{itemize}
    \item Onde o corpo está?
    \item Quanto ele se deslocou?
    \item Quanto tempo levou?
    \item Qual é sua velocidade?
    \item Sua velocidade está aumentando ou diminuindo?
\end{itemize}

Neste estudo, vamos analisar principalmente:

\[
\boxed{
\text{posição}
\longrightarrow
\text{velocidade}
\longrightarrow
\text{aceleração}
}
\]

\end{ideia}


% ==================================================
\section{Velocidade escalar média}
% ==================================================

\begin{ideia}

A \textbf{velocidade escalar média} indica quanto de distância
um corpo percorre em determinado intervalo de tempo.

\[
\boxed{
v_m=\frac{\Delta s}{\Delta t}
}
\]

onde:

\begin{itemize}
    \item $v_m$ é a velocidade escalar média;
    \item $\Delta s$ é a distância percorrida;
    \item $\Delta t$ é o intervalo de tempo.
\end{itemize}

\end{ideia}


\begin{exemplo}

Um carro percorre $180\,\text{km}$ em $3\,\text{h}$.

Aplicando a expressão:

\[
v_m=\frac{\Delta s}{\Delta t}
\]

\[
v_m=\frac{180}{3}
\]

Portanto:

\[
\boxed{
v_m=60\,\text{km/h}
}
\]

Isso significa que, em média, o carro percorreu
$60\,\text{km}$ a cada hora.

\end{exemplo}


\subsection*{Exercícios}

\begin{enumerate}

\item Um ciclista percorre $20\,\text{km}$ em $2\,\text{h}$.
Qual é sua velocidade escalar média?

\vspace{1.5cm}


\item Um carro percorre $150\,\text{km}$ em $2{,}5\,\text{h}$.
Determine sua velocidade escalar média.

\vspace{1.5cm}


\item Um corredor percorre $400\,\text{m}$ em $80\,\text{s}$.
Determine sua velocidade escalar média em $\text{m/s}$.

\vspace{1.5cm}


\item Um automóvel percorre $120\,\text{km}$ nas primeiras
$2\,\text{h}$ de uma viagem e mais $80\,\text{km}$ na hora seguinte.

Determine a velocidade escalar média durante toda a viagem.

\vspace{1.5cm}


\item Um motorista percorre $60\,\text{km}$ a $60\,\text{km/h}$
e depois mais $60\,\text{km}$ a $30\,\text{km/h}$.

Determine a velocidade escalar média de toda a viagem.

\begin{atencao}

Não basta calcular a média aritmética das duas velocidades.

Para encontrar a velocidade média da viagem, devemos calcular:

\[
\boxed{
v_m=
\frac{\text{distância total}}
{\text{tempo total}}
}
\]

\end{atencao}

\end{enumerate}


% ==================================================
\section{Distância e deslocamento}
% ==================================================

\begin{ideia}

A \textbf{distância percorrida} corresponde ao comprimento total
do caminho realizado pelo corpo.

O \textbf{deslocamento} corresponde à diferença entre a posição
final e a posição inicial:

\[
\boxed{
\Delta\vec r=\vec r_f-\vec r_i
}
\]

Portanto, distância e deslocamento são conceitos diferentes.

\end{ideia}


\begin{exemplo}

Um estudante caminha $100\,\text{m}$ para a direita e depois
retorna $30\,\text{m}$.

A distância percorrida é:

\[
d=100+30
\]

\[
\boxed{
d=130\,\text{m}
}
\]

O deslocamento é:

\[
\Delta x=100-30
\]

\[
\boxed{
\Delta x=70\,\text{m}
}
\]

Portanto:

\[
\boxed{
\text{distância}=130\,\text{m}
}
\]

\[
\boxed{
\text{deslocamento}=70\,\text{m}
}
\]

\end{exemplo}


\begin{atencao}

\textbf{Distância:} representa quanto o corpo realmente percorreu.

\textbf{Deslocamento:} representa quanto a posição do corpo mudou.

Um corpo pode percorrer uma grande distância e terminar exatamente
no ponto onde começou.

Nesse caso:

\[
\boxed{
\Delta\vec r=0
}
\]

mesmo que a distância percorrida seja diferente de zero.

\end{atencao}


\subsection*{Exercícios}

\begin{enumerate}

\item Um aluno caminha $50\,\text{m}$ em linha reta para a direita.

Determine:

\begin{enumerate}
    \item a distância percorrida;
    \item o módulo do deslocamento.
\end{enumerate}

\vspace{1.5cm}


\item Uma pessoa caminha $80\,\text{m}$ para a direita e depois
$20\,\text{m}$ para a esquerda.

Determine:

\begin{enumerate}
    \item a distância percorrida;
    \item o deslocamento.
\end{enumerate}

\vspace{1.5cm}


\item Um atleta corre uma volta completa em uma pista circular
de $400\,\text{m}$, retornando exatamente ao ponto de partida.

Determine:

\begin{enumerate}
    \item a distância percorrida;
    \item o deslocamento.
\end{enumerate}

\vspace{1.5cm}


\item Um estudante sai da posição $x=10\,\text{m}$, caminha até
$x=70\,\text{m}$ e depois retorna até $x=30\,\text{m}$.

Determine:

\begin{enumerate}
    \item a distância total percorrida;
    \item o deslocamento.
\end{enumerate}

\vspace{1.5cm}


\item Um móvel realiza o seguinte percurso:

\[
A\rightarrow B=100\,\text{m}
\]

\[
B\rightarrow C=60\,\text{m}
\]

\[
C\rightarrow D=40\,\text{m}
\]

Sabendo que $A$ e $D$ estão separados por $80\,\text{m}$,
determine:

\begin{enumerate}
    \item a distância percorrida;
    \item o módulo do deslocamento;
    \item explique por que os dois valores são diferentes.
\end{enumerate}

\end{enumerate}


% ==================================================
\section{Velocidade vetorial média}
% ==================================================

\begin{ideia}

A \textbf{velocidade vetorial média} utiliza o deslocamento,
e não a distância percorrida.

\[
\boxed{
\vec v_m=
\frac{\Delta\vec r}{\Delta t}
}
\]

Por ser uma grandeza vetorial, ela possui:

\begin{itemize}
    \item módulo;
    \item direção;
    \item sentido.
\end{itemize}

\end{ideia}


\begin{exemplo}

Um aluno desloca-se $100\,\text{m}$ para leste em $20\,\text{s}$.

\[
v_m=\frac{100}{20}
\]

\[
\boxed{
\vec v_m=
5\,\text{m/s para leste}
}
\]

\end{exemplo}


\subsection*{Exercícios}

\begin{enumerate}

\item Um móvel desloca-se $40\,\text{m}$ para a direita em
$10\,\text{s}$.

Determine sua velocidade vetorial média.

\vspace{1.5cm}


\item Um ciclista desloca-se $150\,\text{m}$ para o norte em
$30\,\text{s}$.

Determine o módulo da velocidade vetorial média.

\vspace{1.5cm}


\item Um estudante caminha $100\,\text{m}$ para leste e depois
$40\,\text{m}$ para oeste, realizando o percurso em $20\,\text{s}$.

Determine:

\begin{enumerate}
    \item a distância percorrida;
    \item o deslocamento;
    \item a velocidade vetorial média.
\end{enumerate}

\vspace{1.5cm}


\item Um atleta percorre uma pista circular e retorna ao ponto
de partida após $60\,\text{s}$.

Determine sua velocidade vetorial média.

\vspace{1.5cm}


\item Um barco desloca-se $300\,\text{m}$ para leste e depois
$400\,\text{m}$ para norte em $100\,\text{s}$.

Determine:

\begin{enumerate}
    \item o módulo do deslocamento;
    \item o módulo da velocidade vetorial média;
    \item a direção do deslocamento em relação ao leste.
\end{enumerate}

\end{enumerate}


% ==================================================
\section{Velocidade instantânea}
% ==================================================

\begin{ideia}

A \textbf{velocidade instantânea} representa a velocidade de um
corpo em determinado instante.

É a ideia associada ao valor indicado pelo velocímetro de um carro.

Por exemplo:

\begin{quote}
``Neste instante, o carro está a $80\,\text{km/h}$.''
\end{quote}

Matematicamente:

\[
\boxed{
\vec v=
\frac{d\vec r}{dt}
}
\]

Em uma dimensão:

\[
\boxed{
v=
\frac{dx}{dt}
}
\]

\end{ideia}


\begin{atencao}

Uma maneira simples de pensar é:

\[
\boxed{
\text{velocidade média}
\rightarrow
\text{intervalo de tempo}
}
\]

enquanto:

\[
\boxed{
\text{velocidade instantânea}
\rightarrow
\text{um instante}
}
\]

\end{atencao}


\subsection*{Exercícios}

\begin{enumerate}

\item Um carro está com o velocímetro indicando
$72\,\text{km/h}$.

Qual é sua velocidade instantânea?

\vspace{1.5cm}


\item Converta $72\,\text{km/h}$ para $\text{m/s}$.

\vspace{1.5cm}


\item Um carro apresenta velocidade de $20\,\text{m/s}$ no instante
$t=5\,\text{s}$ e $30\,\text{m/s}$ no instante $t=10\,\text{s}$.

Determine sua aceleração média nesse intervalo.

\vspace{1.5cm}


\item A posição de um móvel é dada por:

\[
x(t)=2t^2+3t
\]

Determine a expressão da velocidade instantânea $v(t)$.

\vspace{1.5cm}


\item A posição de um móvel é dada por:

\[
x(t)=t^3-3t^2+2t
\]

Determine:

\begin{enumerate}
    \item a velocidade instantânea $v(t)$;
    \item a aceleração instantânea $a(t)$;
    \item a velocidade no instante $t=2\,\text{s}$;
    \item a aceleração no instante $t=2\,\text{s}$.
\end{enumerate}

\end{enumerate}


% ==================================================
\section{Aceleração média}
% ==================================================

\begin{ideia}

A \textbf{aceleração} indica como a velocidade de um corpo
varia ao longo do tempo.

A aceleração média é dada por:

\[
\boxed{
a_m=
\frac{\Delta v}{\Delta t}
}
\]

ou:

\[
\boxed{
a_m=
\frac{v_f-v_i}{\Delta t}
}
\]

\end{ideia}


\begin{exemplo}

Um carro aumenta sua velocidade de $10\,\text{m/s}$ para
$30\,\text{m/s}$ em $5\,\text{s}$.

\[
a_m=
\frac{30-10}{5}
\]

\[
a_m=
\frac{20}{5}
\]

Portanto:

\[
\boxed{
a_m=4\,\text{m/s}^2
}
\]

Isso significa que, em média, a velocidade aumenta
$4\,\text{m/s}$ a cada segundo.

\end{exemplo}


\begin{atencao}

A aceleração não significa apenas ``ficar mais rápido''.

A aceleração está relacionada à \textbf{variação do vetor velocidade}.

Portanto, pode existir aceleração quando:

\begin{itemize}
    \item a rapidez aumenta;
    \item a rapidez diminui;
    \item a direção do movimento muda.
\end{itemize}

\end{atencao}


\subsection*{Exercícios}

\begin{enumerate}

\item Um carro aumenta sua velocidade de $10\,\text{m/s}$ para
$20\,\text{m/s}$ em $5\,\text{s}$.

Determine a aceleração média.

\vspace{1.5cm}


\item Uma motocicleta passa de $5\,\text{m/s}$ para $25\,\text{m/s}$
em $4\,\text{s}$.

Determine sua aceleração média.

\vspace{1.5cm}


\item Um carro reduz sua velocidade de $30\,\text{m/s}$ para
$10\,\text{m/s}$ em $5\,\text{s}$.

Determine a aceleração média.

\vspace{1.5cm}


\item Um móvel parte do repouso e atinge $36\,\text{m/s}$ em
$6\,\text{s}$.

Determine:

\begin{enumerate}
    \item a aceleração média;
    \item o significado físico do resultado.
\end{enumerate}

\vspace{1.5cm}


\item Um automóvel aumenta sua velocidade de $10\,\text{m/s}$ para
$25\,\text{m/s}$ durante $3\,\text{s}$. Em seguida, reduz sua velocidade
para $5\,\text{m/s}$ durante mais $4\,\text{s}$.

Determine:

\begin{enumerate}
    \item a aceleração média no primeiro intervalo;
    \item a aceleração média no segundo intervalo;
    \item a aceleração média durante todo o movimento.
\end{enumerate}

\end{enumerate}


% ==================================================
\section{Aceleração instantânea}
% ==================================================

\begin{ideia}

A \textbf{aceleração instantânea} indica como a velocidade está
variando em determinado instante.

Matematicamente:

\[
\boxed{
\vec a=
\frac{d\vec v}{dt}
}
\]

Em uma dimensão:

\[
\boxed{
a=
\frac{dv}{dt}
}
\]

Como:

\[
v=\frac{dx}{dt}
\]

podemos escrever:

\[
\boxed{
a=
\frac{d^2x}{dt^2}
}
\]

Assim, temos a sequência:

\[
\boxed{
x(t)
\longrightarrow
v(t)
\longrightarrow
a(t)
}
\]

\end{ideia}


\subsection*{Exercícios}

\begin{enumerate}

\item Um carro possui velocidade constante de
$20\,\text{m/s}$.

Qual é sua aceleração?

\vspace{1.5cm}


\item Um móvel apresenta velocidade:

\[
v(t)=4t
\]

Determine sua aceleração.

\vspace{1.5cm}


\item Um móvel apresenta velocidade:

\[
v(t)=2t+5
\]

Determine:

\begin{enumerate}
    \item a aceleração;
    \item a velocidade no instante $t=3\,\text{s}$.
\end{enumerate}

\vspace{1.5cm}


\item A posição de um móvel é dada por:

\[
x(t)=2t^2+3t
\]

Determine:

\begin{enumerate}
    \item a velocidade instantânea $v(t)$;
    \item a aceleração instantânea $a(t)$.
\end{enumerate}

\vspace{1.5cm}


\item A posição de um móvel é dada por:

\[
x(t)=t^3-3t^2+2t
\]

Determine:

\begin{enumerate}
    \item a velocidade instantânea $v(t)$;
    \item a aceleração instantânea $a(t)$;
    \item a velocidade no instante $t=2\,\text{s}$;
    \item a aceleração no instante $t=2\,\text{s}$.
\end{enumerate}

\end{enumerate}


% ==================================================
\section{Mapa mental da Cinemática}
% ==================================================

\begin{ideia}

Uma maneira simples de organizar os conceitos fundamentais é:

\[
\boxed{
\text{POSIÇÃO}
}
\]

\[
\downarrow
\]

\[
\boxed{
\text{VELOCIDADE}
}
\]

\[
\downarrow
\]

\[
\boxed{
\text{ACELERAÇÃO}
}
\]

Portanto:

\[
\boxed{
\text{posição}
\longrightarrow
\text{velocidade}
\longrightarrow
\text{aceleração}
}
\]

\end{ideia}


\subsection*{Perguntas para resolver problemas}

Antes de escolher uma fórmula, pergunte:

\begin{enumerate}

\item O problema fala sobre \textbf{distância percorrida}?

\[
\Rightarrow
\text{pense em velocidade escalar média.}
\]

\item O problema fornece \textbf{posição inicial e posição final}?

\[
\Rightarrow
\text{pense em deslocamento.}
\]

\item O problema pergunta quanto a posição varia por unidade de tempo?

\[
\Rightarrow
\text{pense em velocidade.}
\]

\item O problema pergunta quanto a velocidade varia por unidade de tempo?

\[
\Rightarrow
\text{pense em aceleração.}
\]

\item O problema utiliza expressões como
``neste instante''?

\[
\Rightarrow
\text{pense em uma grandeza instantânea.}
\]

\end{enumerate}


% ==================================================
\section{Fórmulas fundamentais}
% ==================================================

\begin{center}

\[
\boxed{
v_m=
\frac{\Delta s}{\Delta t}
}
\]

\vspace{0.3cm}

\[
\boxed{
\vec v_m=
\frac{\Delta\vec r}{\Delta t}
}
\]

\vspace{0.3cm}

\[
\boxed{
\vec a_m=
\frac{\Delta\vec v}{\Delta t}
}
\]

\vspace{0.3cm}

\[
\boxed{
\vec v=
\frac{d\vec r}{dt}
}
\]

\vspace{0.3cm}

\[
\boxed{
\vec a=
\frac{d\vec v}{dt}
}
\]

\vspace{0.3cm}

\[
\boxed{
a=
\frac{d^2x}{dt^2}
}
\]

\end{center}


% ==================================================
% REFERÊNCIA
% ==================================================

\section*{Referência}

A compreensão da Física deve partir da construção de modelos,
da interpretação dos fenômenos e da relação entre conceitos
matemáticos e situações físicas.

A abordagem apresentada busca priorizar a compreensão conceitual
antes da aplicação mecânica das fórmulas
\cite{feynman1964feynman}.


% ==================================================
% BIBLIOGRAFIA
% ==================================================

\printingbibliography

\end{document}